\Needspace{6\baselineskip}
\begin{center}\hypertarget{native-256-article183}{}\pdfbookmark[2]{183.}{article-183}183.\end{center}
The method of the preceding article gives a convenient expression for the measure of precision only for the unknown assigned the last place in the elimination. To avoid this inconvenience, it will be useful to express the coefficient $\delta'''$ in another way. From the equations
\[
\begin{aligned}
P&=p'\\
Q&=q'+\frac{\beta}{\alpha}p'\\
R&=r'+\frac{\gamma'}{\beta'}q'+\frac{\gamma}{\alpha}p'\\
S&=s'+\frac{\delta''}{\gamma''}r'+\frac{\delta'}{\beta'}q'+\frac{\delta}{\alpha}p'
\end{aligned}
\]
it follows that $p'$, $q'$, $r'$, $s'$ can be expressed in terms of $P$, $Q$, $R$, $S$ as follows
\SourceLeaf{257}
\[
\begin{aligned}
p'&=P\\
q'&=Q+\mathfrak{A}P\\
r'&=R+\mathfrak{B}'Q+\mathfrak{A}'P\\
s'&=S+\mathfrak{C}''R+\mathfrak{B}''Q+\mathfrak{A}''P,
\end{aligned}
\]
where $\mathfrak{A}$, $\mathfrak{A}'$, $\mathfrak{B}'$, $\mathfrak{A}''$, $\mathfrak{B}''$, $\mathfrak{C}''$ are determinate quantities. Hence we shall have (restricting the number of unknowns to four)
\[
s=-\frac{\lambda'''}{\delta'''}+\frac{\mathfrak{A}''}{\delta'''}P
+\frac{\mathfrak{B}''}{\delta'''}Q+\frac{\mathfrak{C}''}{\delta'''}R
+\frac{1}{\delta'''}S.
\]
We draw the following conclusion. The most probable values of the unknowns $p$, $q$, $r$, $s$, etc. are to be derived by elimination from the equations $P=0$, $Q=0$, $R=0$, $S=0$, etc. If we regard $P$, $Q$, $R$, $S$, etc. for the moment as undetermined, the same elimination procedure will clearly express the unknowns in linear form in terms of $P$, $Q$, $R$, $S$, etc., giving
\[
\begin{gathered}
\begin{aligned}
p&=L+AP+BQ+CR+DS+\text{etc.}\\
q&=L'+A'P+B'Q+C'R+D'S+\text{etc.}\\
r&=L''+A''P+B''Q+C''R+D''S+\text{etc.}\\
s&=L'''+A'''P+B'''Q+C'''R+D'''S+\text{etc.}
\end{aligned}\\
\text{etc.}
\end{gathered}
\]
Once this is done, the most probable values of $p$, $q$, $r$, $s$, etc. will clearly be $L$, $L'$, $L''$, $L'''$, etc., respectively. The measures of precision to be assigned to these determinations will be expressed respectively by
$\BareRoot{\dfrac{1}{A}}\mathpunct{\raisebox{.45ex}{,}}$
$\BareRoot{\dfrac{1}{B'}}\mathpunct{\raisebox{.45ex}{,}}$
$\BareRoot{\dfrac{1}{C''}}\mathpunct{\raisebox{.45ex}{,}}$
$\BareRoot{\dfrac{1}{D'''}}$, etc., with the precision of the original observations set $=1$. What we have already proved about the determination of the unknown $s$ (for which $\dfrac{1}{\delta'''}$ corresponds to $D'''$) can be extended to all the others simply by permuting the unknowns.%
\HistoricalUnmarkedNote{[Manuscript note:] If the degrees of precision of the quantities $a$, $b$, $c$, etc. are $\alpha$, $\beta$, $\gamma$, etc., the degree of precision of the quantity $x=a+b+c+\text{etc.}$ will be $\dfrac{1}{\BareRoot{\left(\dfrac{1}{\alpha\alpha}+\dfrac{1}{\beta\beta}+\dfrac{1}{\gamma\gamma}+\text{etc.}\right)}}\mathpunct{\cdot}$}
